% !TeX root = ../Tex/main.tex

\section{Literature Review and Limitations}

\subsection{A Review of Models of Conflicts}
Further exploration of the literature on games of conflict may support as well as weaken the analysis presented so far. In any case, already-established models are a wonderful comparison. Some of the papers that stand out are \parencite{schellingStrategyConflict1960, Schelling1966}. On the other hand, I think it may be of interest to focus on \textcite{mckelveyQuantalResponseEquilibria1995} to further develop the model.\

\subsubsection{A Review of Models of Conflicts with a Focus on Bargaining}
A specific branch of the literature focuses on bargaining models of conflicts. Among the most relevant contributions to this fiels are from \textcite{Fearon1996, Fearon_1995, Fearon_1998} and again \textcite{schellingStrategyConflict1960, Schelling1966}.
\textcite{reiterExploringBargainingModel2003} offers an in depth review of the field. He provides an historical perspective as well as recent contributions. Additionally, he sums up the most important ideas.

\vspace*{1em}
\textcite{reiterExploringBargainingModel2003} says: "James Fearon prominently contributed to this scholarship with his 1995 paper "Rationalist Explanations for War." He usefully highlighted the point that if states agreed on the outcome of a possible war, they could probably avoid war. Central to Fearon’s paper is the importance of focusing on states' disagreement over their capabilities and/or resolve. [\dots] Fearon uses a bargaining model to argue that if two states in dispute know the outcome of a possible war, they should in general prefer to reach a deal that would reflect the hypothetical postwar political settlement, rather than fight, reach that same settlement, and also suffer the costs of war. This view assumes that fighting itself is costly—that the belligerent always suffers some negative utility, no matter how the issue at stake is settled. In fact, this assumption becomes necessary when one asserts that two states would always prefer to reach a bargain without fighting rather than fight and then reach the same bargain. Within his bargaining model, Fearon develops three conditions under which war is possible."

\vspace*{1em}
\textcolor{purple}{Modelli di cherim:} In order to broaden the overview, some other interesting contributions come from
\textcite{Kirshner01092000, wagnerBargainingWar2000, wittmanHowWarEnds1979, smithBargainingNatureWar2004, filsonBargainingFightingImpact2004}.

\subsection{Limits of Theoretical Models}
In many courses I covered the topic of formal modelling. In the previous pages I present my personal reelaboration along with a brief summary of all litarature I studied on this topic. This is for sure not enough to assess the state of models of this article. However, the exposition I provide the reader with, should be considered satisfactory in the sense that there will never be any such thing as \textit{the correct, perfect or true modelling theory.}

Under this respect, it is possible to further expand the discussion on limitations. The literature offers some very good papers, among others, I report:
\textcite{ashworthTheoryCredibilityIntegrating2021, brauningerTheoryBuildingCausal2020, alma991001906059707991}. The present different ways to combine empirical and theoretical researc along with some very interersting interpretations of the meaning and scope of models.

\subsection{Connection between Game Theory and Biology}
There are many great applications of game theory to biology. I believe the most interesting ones are: \textcite{smithLogicAnimalConflict1973, axelrodEvolutionCooperation1981}.
Future works may shed light on more recent, as well as interesting applications.
A vaste literature covers the topic of evolutionary games too: \textcite{kandoriLearningMutationLong1993, taylorEvolutionaryStableStrategies1978, hofbauerEvolutionaryGamesPopulation2003, nowakFiveRulesEvolution2006}. In particular, \textcite[chap 12]{hofbauerEvolutionaryGamesPopulation2003} offers a clear treaty of powerful analytical tools along with applications to these game theoretic tools to biology.

Unfortunately, an extension dedicated to these numerous works must be postponed.